\documentclass[11pt]{amsart}
\usepackage{geometry}                % See geometry.pdf to learn the layout options. There are lots.
\geometry{letterpaper}                   % ... or a4paper or a5paper or ... 
%\geometry{landscape}                % Activate for for rotated page geometry
%\usepackage[parfill]{parskip}    % Activate to begin paragraphs with an empty line rather than an indent
\usepackage{graphicx}
\usepackage{amssymb}
\usepackage{epstopdf}
\DeclareGraphicsRule{.tif}{png}{.png}{`convert #1 `dirname #1`/`basename #1 .tif`.png}

\title{PS 2.6}
%\date{}                                           % Activate to display a given date or no date

\begin{document}
\maketitle
%\section{}
%\subsection{}
 \noindent  Haga un diagrama de flujo para calcular el precio del billete ida y vuelta en ferrocarril, conociendo la distancia del viaje de ida y el tiempo de estancia. Se sabe además que si el número de días de estancia es superior a 7 y la distancia total (ida y vuelta) a recorrer es superior a 800 km, el billete tiene una reducción del 30\%. El precio por km es de \$0.23. \newline

 \noindent Datos: \hspace{0.5cm} DIST, TIEM.
 
 \noindent Donde: 
 
 \setlength{\parindent}{1cm}
 DIST \hspace{0.5cm}  es una variable de tipo entero que representa la distancia del viaje de ida. 
 
 \setlength{\parindent}{1cm}
 TIEM \hspace{0.5cm} es una variable de tipo entero que representa el tiempo de estancia.
\end{document}  